% ════════════════════════════════════════════════════════════════════════════
\section{Notation, and two refinements of the model}
\label{app:notation}
% ════════════════════════════════════════════════════════════════════════════

\subsection{Notation}

\noindent The symbols below recur throughout; each is defined in full where it is
introduced, and all are gathered here for reference.

\vspace{4pt}

\begingroup\small
\renewcommand{\arraystretch}{1.2}
\begin{tabularx}{\linewidth}{@{}l X@{}}
\toprule
\multicolumn{2}{@{}l}{\textbf{Principals, assets, economics}}\\
\midrule
$\Victim,\ \Attacker$ & holder (victim) and coercer (attacker) \\
$K$ & master secret gating the asset \\
$W$ & value of the asset to $\Victim$ \\
$c$ & expected disutility to $\Attacker$ of killing $\Victim$ \\
$\pi$ & $\Victim$'s \emph{pivotality} --- win-probability $\Attacker$ gains by removing $\Victim$ \\
\addlinespace
\multicolumn{2}{@{}l}{\textbf{Adversary, game, safety}}\\
\midrule
WDY & the Wrench--Dolev--Yao adversary, fixed by clauses \ref{d:kerckhoffs}--\ref{d:nodelegatecoerce} (Assumption~\ref{ass:power}) \\
$p_A$ & probability $\Attacker$ reaches $K$ after the encounter; \emph{gray} iff $0<p_A<1$ \\
$t$ & coercion-time budget (how long $\Attacker$ can sustain coercion) \\
$G$ & protocol time-gate (enforced latency to complete the derivation) \\
$\mathsf{Corr}$ & residual-window event: the points $\Victim$ has surrendered so far are correct (\S\ref{sec:residual-dr}) \\
$\mathsf{Esc}$ & residual-window event: the ${\approx}1$--$2$ min shaved by brute-forcing the last points decides escape from the scene (\S\ref{sec:residual-dr}) \\
\addlinespace
\multicolumn{2}{@{}l}{\textbf{Construction (TGPO orbit \S\ref{sec:gw})}}\\
\midrule
$\sigma$ & public, precomputation-resistant timechain salt\codename{Namtso} \\
$\Hcheap,\ \Hhard$ & unspecified cheap hash; unspecified memory- and depth-hard (Argon2-class) hash \\
$o_i$ & orbit state at stage $i \in \{0,\dots,N\}$; $o_i = \Hhard(u_{i-1})$ (with $o_0 = \Hcheap(\sigma)$); $N{+}1$ stages total \\
$N$ & number of \emph{deep} stages --- those requiring at least one memory- and depth-hard hash ($\Hhard$) to derive \\
$\Sh_i$ & Shamir polynomial at stage $i$; $r_i$ of its points reconstruct it \\
$u_i$ & $\Hcheap(o_{i}\|\Sh_{i})$: commitment of the wiped values $o_{i}$ and $\Sh_{i}$ \\
$K_i$ & $\Hcheap(o_{i}\|\Sh_{i})$: master secret correspondent to stage $i$. \\
$p_{i,j},\ \theta_{i,j}$ & point $j$ of stage $i$; parameters for fractal $j$ of stage $i$ \\
$r_i$ & Shamir threshold at stage $i$ ($r_i \ge 2$) \\
$\sftstar$ & $2^{64}$ memory- and depth-hard evaluations (infeasible) \\
floor & \emph{prohibitive} cost: $\ge \sftstar$ hard \emph{or} $\ge 2^{96}$ cheap (Stipulation~\ref{ass:feas}) \\
\bottomrule
\end{tabularx}
\endgroup

\subsection{The two tacit-seizure clauses}
\label{app:notation-tacit}

\begin{remark}[The two non-seizure clauses coincide for design]\label{rem:tacit-coincide}
For protocol-design purposes the tacitness of the entropy-carrying secrets
(\ref{d:notacit}) and that of the perceptual-oracle parameters
(\ref{d:noperceptual}) are logically equivalent. If the parameters were
transmissible, a wrench attack would seize them --- and, available within that
\emph{same} encounter, they would then serve to seize the entropy-carrying points
themselves; so non-tacit parameters force non-tacit points. Conversely, if the
points were somehow transmissible without the parameters ever being available,
then the tacitness of the parameters is irrelevant, the points having already
fallen. Either way the binary question ``does a single seizure reach $K$?'' turns
on the two clauses together, and the model may treat them as one.
\end{remark}

\begin{remark}[Where the distinction still matters]\label{rem:tacit-quantitative}
The distinction can nonetheless matter for a \emph{quantitative} refinement --- a
statistical distribution of the wall-clock needed to coerce either (up to sufficient
number of most-significant bits) given $\Attacker$'s coercion time budget $t$ 
(Definition \ref{def:priced}). Still, a secret that is \emph{theoretically} transmissible, 
but only with pen and paper, an adversarial fork of the application, and long hours 
to days of focused mental effort under perfect $\Attacker$--$\Victim$ cooperation, 
is essentially as good as ideal. Until empirical studies of TKBA (Assumption~\ref{ass:tkba}) 
are done, the actual relevance of that nuance is purely speculative.
\end{remark}

\begin{remark}[TKBA restates the tacit-seizure clauses]\label{rem:tkba-restates}
Assumption~\ref{ass:tkba} is not an independent hypothesis: it is
clauses~\ref{d:notacit}--\ref{d:noperceptual} lifted out of the adversary model and
stated as a premise on the secret. The clauses say $\Attacker$ cannot seize the
tacit points or the perceptual parameters even under full cooperation
(\ref{d:subjugation}); TKBA says the same thing about the secret itself ---
recognizable but not transmissible --- with no adversary mentioned. They are two
faces of one assumption (exactly as \ref{d:kerckhoffs} and
Assumption~\ref{ass:kerckhoffs} are two faces of Kerckhoffs), and by
Remark~\ref{rem:tacit-coincide} the two clauses are themselves one. We keep both
forms because each is the natural one in its place --- the premise where the
construction is motivated (\S\ref{sec:gw}), the adversary clause where seizure is
reasoned about.
\end{remark}

